% Generated by audit_reproduce.py from results_3d.csv; do not hand-edit numbers.
В скрученном кубе адаптивный вариант при $\mu=0.1$ даёт
$CV_V=0.0096$ против $0.0239$ у inverse mean ratio
(снижение на $59.7\,\%$). Значения ортогональности близки, но не тождественны:
$0.985550$ и $0.985561$ соответственно.
В шаре inverse mean ratio имеет наибольший среди трёх методов
$Q_{\rm orth}=0.9897$ и наименьший
$AR_{95}=2.1021$; адаптивный вариант даёт меньший разброс
объёмов: $CV_V=0.2212$ против $0.3684$,
но уступает по ортогональности.

В арковой призме inverse mean ratio имеет наибольший $Q_{\rm orth}$ и
наибольшее значение минимального углового масштабированного якобиана, а также наименьшие $CV_V$ и $AR_{95}$
в проверенной тройке. Его $AR_{95}=4.9563$, тогда как для
адаптивного варианта $AR_{95}=8.2422$.
Однако по размерному показателю $\sigma_{\rm len}$ адаптивный вариант лучше:
$0.1143$ против $0.1154$.
Поэтому превосходство по \emph{всем} метрикам не утверждается.
Это наблюдения при указанных разрешениях, границе, начальной сетке и параметрах,
а не универсальный рейтинг функционалов и не оценка ошибки решения уравнения.
